\documentclass[11pt,a4paper]{article}

\usepackage{kotex}
\usepackage{fontspec}
\setmainfont{Times New Roman}
\setmainhangulfont{AppleMyungjo}
\setsanshangulfont{Apple SD Gothic Neo}

\usepackage[a4paper,top=18mm,bottom=18mm,left=20mm,right=20mm]{geometry}
\usepackage{fancyhdr}
\setlength{\headheight}{24pt}
\usepackage{enumitem}
\usepackage{mdframed}
\usepackage{array}
\usepackage{booktabs}

\setlength{\parindent}{1em}
\linespread{1.18}

\pagestyle{fancy}
\fancyhf{}
\renewcommand{\headrulewidth}{0.6pt}
\fancyhead[L]{\sffamily\bfseries 2026 고1 영어 변형 — 지문 32}
\fancyhead[R]{\sffamily [다양한 장소에서 공부하기]}
\renewcommand{\footrulewidth}{0pt}
\fancyfoot[C]{\framebox[16mm][c]{\thepage}}

\newcommand{\qno}[1]{\par\vspace{8pt}\noindent{\sffamily\bfseries #1.}\ }
\newcommand{\instr}[1]{{\sffamily #1}\par\vspace{3pt}}
\newlist{choices}{enumerate}{1}
\setlist[choices]{label=,leftmargin=1.5em,itemsep=2pt,topsep=3pt,parsep=0pt}
\newmdenv[linewidth=0.6pt,roundcorner=3pt,innertopmargin=7pt,innerbottommargin=7pt,
          innerleftmargin=9pt,innerrightmargin=9pt,skipabove=5pt,skipbelow=5pt]{pbox}
\newcommand{\nemo}[1]{\fbox{\,#1\,}}

\begin{document}

\begin{center}
{\fontsize{15}{17}\selectfont\sffamily\bfseries [지문 32] 다음 글을 읽고 물음에 답하시오.}
\end{center}
\vspace{4pt}

\begin{pbox}
Do you always study in the same spot? Don't. If you study in different places,
that helps create variety, and that rich experience can \rule{3cm}{0.4pt} what
you are learning. When several college students told us they did school work in
different places rather than in some favorite nook, that habit reflected the
research on learning. Numerous experiments have found that if learners simply
study in at least two different places, they are more likely to recall the
material. In one of the first such trials, two different groups studied a list
of words. Some students returned to the same room twice while their counterparts
spent the same amount of time divided between two locations. When asked to recall
as many words on the list as possible, those who had moved around did far better.
Variety creates rich association, even when those connections form in the
background, totally outside of what we are consciously thinking.

\medskip
\footnotesize
* nook: 구석진 곳\quad ** counterpart: 상대
\end{pbox}

\bigskip

\qno{1}\instr{다음 빈칸에 들어갈 말로 가장 적절한 것은?}
\begin{choices}
\item ① replace
\item ② complicate
\item ③ contradict
\item ④ consolidate
\item ⑤ accelerate
\end{choices}

\qno{2}\instr{윗글의 제목으로 가장 적절한 것은?}
\begin{choices}
\item ① Why Comfortable Study Environments Improve Concentration
\item ② Study in Multiple Places to Remember More
\item ③ The Science Behind Why Students Forget What They Learn
\item ④ How Group Study Outperforms Solo Learning
\item ⑤ Creating a Dedicated Study Space for Better Habits
\end{choices}

\qno{3}\instr{[서술형] 다음 문장을 우리말로 해석하시오. \normalfont(2점)}

\begin{quote}
\itshape
Variety creates rich association, even when those connections form in the
background, totally outside of what we are consciously thinking.
\end{quote}

\noindent $\Rightarrow$ \rule{\linewidth}{0.4pt}

\vspace{4pt}
\noindent \rule{\linewidth}{0.4pt}

% ===================== Q4 =====================
\qno{4}\instr{윗글의 내용과 일치하지 \underline{않는} 것은?}
\begin{choices}
\item ① Studying in different places can create richer learning experiences.
\item ② Students who used two locations recalled more words than those who stayed in one room.
\item ③ The benefit of varied places may arise outside conscious thought.
\item ④ The passage recommends always returning to one favorite study nook.
\item ⑤ Variety can help form associations that strengthen learning.
\end{choices}

\end{document}
